% ==================================================
% Summation Operator
% Discrete Calculus Notes Stub
% ==================================================

\section{Discrete Antiderivatives}

The telescoping identity established in the previous section reveals a deep
relationship between the difference operator and summation.  In particular,

\[
\sum_{k=a}^{b-1} \Delta f(k) = f(b) - f(a).
\]

This identity suggests that summation acts as an inverse operation to the
difference operator, much like integration acts as an inverse to
differentiation in ordinary calculus.

\subsection*{Definition}

\begin{definition}[Discrete Antiderivative]
Let \(f : \mathbb{Z} \to \mathbb{R}\) be a discrete function.
A function \(F : \mathbb{Z} \to \mathbb{R}\) is called a
\emph{discrete antiderivative} of \(f\) if

\[
\Delta F(n) = f(n).
\]

That is,

\[
F(n+1) - F(n) = f(n).
\]
\end{definition}

\begin{remark}
A discrete antiderivative plays the same role in discrete calculus that an
antiderivative plays in ordinary calculus.  Instead of reversing
differentiation, it reverses the forward difference operator.
\end{remark}

\subsection*{Example}

Consider the constant function

\[
f(n) = 1.
\]

We seek a function \(F\) such that

\[
\Delta F(n) = 1.
\]

Taking \(F(n) = n\), we compute

\[
\Delta F(n) = F(n+1) - F(n) = (n+1) - n = 1.
\]

Thus \(F(n) = n\) is a discrete antiderivative of \(1\).

\subsection*{Another Example}

Let

\[
f(n) = n.
\]

One can verify that

\[
F(n) = \frac{n(n-1)}{2}
\]

is a discrete antiderivative since

\[
\Delta F(n)
=
\frac{(n+1)n}{2} - \frac{n(n-1)}{2}
=
n.
\]

\subsection*{Constant of Integration}

Discrete antiderivatives are not unique.

\begin{proposition}
If \(F\) is a discrete antiderivative of \(f\), then so is \(F + C\)
for any constant \(C\).
\end{proposition}

\begin{proof}
\[
\Delta(F+C)(n)
=
(F+C)(n+1) - (F+C)(n)
=
F(n+1)-F(n)
=
\Delta F(n)
=
f(n).
\]
\end{proof}

